\begin{table}[htbp]
\centering
\caption{Standardized Effect Sizes}
\label{tab:sde}
\begin{adjustbox}{max width=\textwidth}
\begin{tabular}{lcccccc}
\hline\hline
Outcome & $\hat{\beta}$ & SE & SD($Y$) & SDE & SE(SDE) & Classification \\
\hline
\textit{Panel A: Pooled} & & & & & & \\
Log total votes & $-0.0141$ & $0.0058$ & $1.5021$ & $-0.0094$ & $0.0039$ & Small negative \\
Turnout rate & $-0.0057$ & $0.0032$ & $0.1155$ & $-0.0493$ & $0.0278$ & Small negative \\
[3pt]
\textit{Panel B: Heterogeneous (sample splits)} & & & & & & \\
Urban counties & $-0.0251$ & $0.0067$ & $1.5021$ & $-0.0167$ & $0.0045$ & Small negative \\
Rural counties & $-0.0178$ & $0.0087$ & $1.5021$ & $-0.0118$ & $0.0058$ & Small negative \\
\hline\hline
\end{tabular}
\end{adjustbox}
\begin{tablenotes}
\small
\item \textit{Notes:} \textbf{Country:} United States. \textbf{Research question:} Does the loss of USPS post office establishments reduce voter turnout in presidential elections? \textbf{Policy mechanism:} The USPS Retail Access Optimization Initiative (RAOI, 2011--2017) closed approximately 480 post offices nationwide as part of a cost-reduction strategy driven by the 2006 Postal Accountability and Enhancement Act's pre-funding mandate, reducing physical postal access points in affected communities. \textbf{Outcome definition:} Log total presidential votes at the county level from MIT Election Data + Science Lab county presidential returns; turnout rate defined as total votes divided by voting-age population proxy (Census ACS population $\times$ 0.76). \textbf{Treatment:} Binary indicator for counties that lost at least one USPS establishment (NAICS 491110) between 2014 and 2018 as measured by the BLS Quarterly Census of Employment and Wages. \textbf{Data:} BLS QCEW (2014--2023), MIT Election Lab county presidential returns (2000--2024), Census ACS 5-year estimates (2015, 2019, 2022); balanced panel of 1,791 counties observed across seven presidential elections. \textbf{Method:} Callaway and Sant'Anna (2021) doubly robust staggered difference-in-differences with never-treated controls; standard errors are analytical. Note: parallel trends pre-test rejects (p = 0.004), so estimates may reflect pre-existing convergence rather than a causal effect. \textbf{Sample:} Counties with non-missing QCEW USPS data across all years (2014--2023) and non-missing presidential vote totals in all seven election cycles; 514 treated, 1,277 never-treated. SDE $= \hat{\beta} / \text{SD}(Y)$ where SD($Y$) is the pre-treatment standard deviation. Classification refers to magnitude, not statistical significance: Large ($|$SDE$| > 0.15$), Moderate ($0.05$--$0.15$), Small ($0.005$--$0.05$), Null ($< 0.005$).
\end{tablenotes}
\end{table}
